% Luc Destombe --- licence CC BY-NC-SA 4.0
% https://creativecommons.org/licenses/by-nc-sa/4.0/deed.fr
% Fichier autonome : compiler avec pdflatex, deux fois (signets, nombre de pages).
\documentclass[11pt,a4paper]{article}
\makeatletter
\newif\ifeleve
\elevefalse

\newif\iflycee@palatino
\lycee@palatinotrue

\RequirePackage[utf8]{inputenc}
\RequirePackage[T1]{fontenc}
\RequirePackage[french]{babel}
\RequirePackage{etoolbox}

\newcommand{\ShorthandsOff}{\@ifundefined{shorthandoff}{}{\shorthandoff{;:!?}}}
\newcommand{\ShorthandsOn}{\@ifundefined{shorthandon}{}{\shorthandon{;:!?}}}
\ShorthandsOff
\AtBeginDocument{\ShorthandsOn}
\AtBeginEnvironment{tikzpicture}{\ShorthandsOff}

\RequirePackage{geometry}
\RequirePackage{fancyhdr}
\RequirePackage{lastpage}
\RequirePackage{multicol}
\RequirePackage{array}
\RequirePackage{enumitem}
\RequirePackage{xcolor}
\RequirePackage{needspace}

\RequirePackage{amsmath,amssymb}
\RequirePackage{mathpazo}
\DeclareMathSizes{10.95}{10.95}{8}{5.5}
\begingroup\catcode`\_=\active \catcode`\^=\active
\gdef_#1{\sb{\scriptscriptstyle #1}}
\gdef^#1{\sp{\scriptscriptstyle\lycee@moinsepais #1}}
\endgroup
\newcommand\lycee@moins{\mathbin{%
  \setbox\z@\hbox{$\m@th\scriptscriptstyle\lycee@moinsorig$}%
  \dimen@=.55\wd\z@ \advance\dimen@-.5pt
  \hbox to.75\wd\z@{\hss$\m@th\scriptscriptstyle\vcenter{\hbox{%
    \tikz\draw[line width=.5pt,line cap=round](0,0)--(\the\dimen@,0);}}$\hss}}}
\begingroup\lccode`\~=`\- \lowercase{\endgroup\let~}\lycee@moins
\newcommand\lycee@moinsepais{\mathcode`\-="8000 }
\AtBeginDocument{\mathchardef\lycee@moinsorig=\mathcode`\-\relax}
\AtBeginDocument{\catcode`\_=12 \mathcode`\_="8000
  \catcode`\^=12 \mathcode`\^="8000 }
\newcommand\lycee@exposants{%
  \fontdimen13\textfont2=.48\fontdimen6\textfont2
  \fontdimen14\textfont2=.48\fontdimen6\textfont2
  \fontdimen15\textfont2=.48\fontdimen6\textfont2
  \fontdimen13\scriptfont2=.48\fontdimen6\scriptfont2
  \fontdimen14\scriptfont2=.48\fontdimen6\scriptfont2
  \fontdimen15\scriptfont2=.48\fontdimen6\scriptfont2}
\every@math@size\expandafter{\the\every@math@size\lycee@exposants}
\newlength{\retraitcalculs}
\setlength{\retraitcalculs}{2em}

\RequirePackage{tikz}
\usetikzlibrary{arrows.meta,calc,babel,shapes.misc}
\RequirePackage[most]{tcolorbox}
\tcbuselibrary{skins,breakable}

\definecolor{coulDef}{HTML}{1F4E79}
\definecolor{coulProp}{HTML}{7030A0}
\definecolor{coulRem}{HTML}{C55A11}
\definecolor{coulEx}{HTML}{2E7D32}
\definecolor{coulExo}{HTML}{B03A2E}
\definecolor{coulPart}{HTML}{1F4E79}
\definecolor{coulMeth}{HTML}{1B6E4F}
\definecolor{coulCourbe}{HTML}{0B5ED7}
\definecolor{coulCourbeB}{HTML}{C00000}
\definecolor{coulCourbeC}{HTML}{2E7D32}
\definecolor{coulGrille}{HTML}{8C8C8C}
\definecolor{coulRep}{HTML}{1B6E4F}
\definecolor{coulBar}{HTML}{2E7D32}

\newcommand{\R}{\mathbb{R}}
\newcommand{\Rs}{\mathbb{R}^{*}}
\renewcommand{\le}{\leqslant}
\renewcommand{\ge}{\geqslant}
\newcommand{\vect}[1]{\overrightarrow{#1}}
\newcommand{\norme}[1]{\left\lVert #1 \right\rVert}
\newcommand{\intervalle}[2]{\left[#1\,;#2\right]}
\RequirePackage{cancel}
\renewcommand{\CancelColor}{\color{coulExo}}
\newcommand{\fois}[1]{\times{\color{coulExo}#1}}
\newcommand{\barre}[1]{\cancel{#1}}

\tikzset{grille/.style={coulGrille,line width=0.4pt}}
\newcommand{\quadrillage}[4]{\draw[grille,step=1] (#1,#2) grid (#3,#4);}
\tikzset{
  croix/.style={cross out,draw,line width=0.6pt,inner sep=0pt,minimum size=4pt},
  croixdroite/.style={croix,rotate=45,line width=0.8pt},
}
\newcommand{\pt}[3]{\path (#1) node[croix]{} node[#2,font=\small,inner sep=2.5pt]{$#3$};}

\newcommand{\lycee@repere}[4]{%
  \draw[grille,step=1] (#1,#2) grid (#3,#4);
  \draw[-{Stealth[length=2mm]},thick] (#1,0)--({#3+0.4},0) node[below left=-1pt]{$x$};
  \draw[-{Stealth[length=2mm]},thick] (0,#2)--(0,{#4+0.4}) node[below left=-1pt]{$y$};
  \node[below left,font=\scriptsize,inner sep=1.5pt] at (0,0) {$0$};}
\newcommand{\lycee@gradx}[1]{\foreach \k/\l in {#1}{%
  \draw (\k,0.07)--(\k,-0.07) node[below,font=\scriptsize,inner sep=1.5pt]{$\l$};}}
\newcommand{\lycee@grady}[1]{\foreach \k/\l in {#1}{%
  \draw (0.07,\k)--(-0.07,\k) node[left,font=\scriptsize,inner sep=1.5pt]{$\l$};}}
\newcommand{\lycee@intff}[2]{\left[#1\,;#2\right]}
\newcommand{\lycee@intfo}[2]{\left[#1\,;#2\right[}
\newcommand{\lycee@intof}[2]{\left]#1\,;#2\right]}
\newcommand{\lycee@intoo}[2]{\left]#1\,;#2\right[}
\newcommand{\lycee@ens}[1]{\left\{#1\right\}}
\AtBeginDocument{%
  \providecommand{\repere}{\lycee@repere}%
  \providecommand{\gradx}{\lycee@gradx}%
  \providecommand{\grady}{\lycee@grady}%
  \providecommand{\Cf}{\mathcal{C}\sb{\scriptscriptstyle f}}%
  \providecommand{\Cg}{\mathcal{C}\sb{\scriptscriptstyle g}}%
  \providecommand{\intff}{\lycee@intff}%
  \providecommand{\intfo}{\lycee@intfo}%
  \providecommand{\intof}{\lycee@intof}%
  \providecommand{\intoo}{\lycee@intoo}%
  \providecommand{\ens}{\lycee@ens}}

\newlist{questions}{enumerate}{3}
\setlist[questions,1]{label=\arabic*\textdegree),leftmargin=*,itemsep=2pt,topsep=3pt}
\setlist[questions,2]{label=\alph*),leftmargin=*,itemsep=1pt,topsep=2pt}
\setlist[questions,3]{label=\roman*),leftmargin=*,itemsep=1pt,topsep=1pt}
\setlist[itemize]{label=\textbullet,leftmargin=*,itemsep=2pt,topsep=3pt}

\newcommand{\besoinplace}[1]{\par\penalty\@M\begingroup
  \dimen@=#1\relax\dimen@ii=\pagegoal\advance\dimen@ii-\pagetotal
  \ifdim\dimen@>\dimen@ii\ifdim\dimen@ii>\z@\vfil\fi\break\fi\endgroup}

\newcommand{\lycee@pied}{\thepage/\pageref{LastPage}}
\newcommand{\entete}[2]{%
  \pagestyle{fancy}\fancyhf{}%
  \fancyhead[L]{\footnotesize\color{coulPart}\textbf{#1}}%
  \fancyhead[R]{\footnotesize\color{coulPart}\textbf{#2}}%
  \fancyfoot[C]{\footnotesize\lycee@pied}%
  \renewcommand{\headrulewidth}{0.4pt}%
  \renewcommand{\headrule}{\hbox to\headwidth{%
    \color{coulPart!40}\leaders\hrule height \headrulewidth\hfill}}}

\RequirePackage{ccicons}
\newcommand{\lycee@licence}{{\scriptsize\color{gray}%
  \copyright\ Luc Destombe --- \ccbyncsa\ CC BY-NC-SA 4.0}}
\newif\iflycee@licence \lycee@licencetrue
\newcommand{\sanslicence}{\lycee@licencefalse}
\AtBeginDocument{\iflycee@licence\AddToHookNext{shipout/foreground}{%
  \put(0,0){\rlap{\kern\dimexpr1in+\hoffset+\oddsidemargin+\textwidth\relax
    \llap{\raisebox{-\dimexpr1in+\voffset+\topmargin+\headheight+\headsep
      +\textheight+\footskip\relax}{\lycee@licence}}}}}\fi}

\newcommand{\formulesserrees}{%
  \setlength{\abovedisplayskip}{3pt}\setlength{\belowdisplayskip}{3pt}%
  \setlength{\abovedisplayshortskip}{2pt}\setlength{\belowdisplayshortskip}{3pt}}

\setlength{\parindent}{0pt}

\RequirePackage[implicit=false,hidelinks,pdfencoding=auto,psdextra,bookmarksopen,
  bookmarksopenlevel=1,pdfcreator={},pdfproducer={}]{hyperref}
\RequirePackage{bookmark}
\hypersetup{pdfauthor={Luc Destombe},pdfkeywords={CC BY-NC-SA 4.0}}
\newcounter{lycee@signet}
\newcommand{\signet}[3][]{\stepcounter{lycee@signet}%
  \edef\lycee@dest{\if\relax\detokenize{#1}\relax
    signet.\arabic{lycee@signet}\else#1\fi}%
  \hypertarget{\lycee@dest}{}\bookmark[level=#2,dest=\lycee@dest]{#3}}
\pdfstringdefDisableCommands{%
  \def\R{R}\def\vect#1{#1}\def\Cf{Cf}\def\Cg{Cg}\def\fois#1{×#1}%
  \def\barre#1{#1}\def\intervalle#1#2{[#1;#2]}\def\intff#1#2{[#1;#2]}%
  \def\textsuperscript#1{#1}\def\dfrac#1#2{#1/#2}\def\frac#1#2{#1/#2}%
  \def\vec#1{#1⃗}}%
\geometry{top=2cm,bottom=1cm,left=1cm,right=1cm,headheight=14pt,footskip=0.6cm}

\RequirePackage{pgfplots}
\pgfplotsset{compat=1.18}
\RequirePackage{tkz-tab}

\newcounter{partiecours}
\renewcommand{\thepartiecours}{\Roman{partiecours}}
\newcommand{\partiecours}[1]{%
  \refstepcounter{partiecours}%
  \Needspace*{8\baselineskip}%
  \bigskip
  \noindent\begin{tcolorbox}[enhanced,colback=coulPart,colframe=coulPart,
      boxrule=0pt,arc=2pt,left=8pt,right=8pt,top=4pt,bottom=4pt,
      before skip=6pt,after skip=10pt]
    \color{white}\bfseries\large \thepartiecours{} --- #1%
    \signet[partie-\arabic{partiecours}]{1}{\thepartiecours{} --- #1}
  \end{tcolorbox}%
  \addcontentsline{toc}{section}{\thepartiecours{} --- #1}%
}

\newtcolorbox{boitecours}[3][]{%
  enhanced,breakable,
  colback=white,colframe=#2,coltitle=white,
  fonttitle=\bfseries,
  attach boxed title to top left={xshift=8pt,yshift=-\tcboxedtitleheight/2},
  boxed title style={colback=#2,arc=2pt,boxrule=0pt},
  boxrule=0.9pt,arc=2pt,
  left=8pt,right=8pt,top=8pt,bottom=6pt,
  title={#3},#1}

\newcounter{methode}
\newenvironment{definitionbox}[1][Définition]%
  {\begin{boitecours}[phantom={\signet{2}{#1}}]{coulDef}{#1}}{\end{boitecours}}
\newenvironment{proprietebox}[1][Propriété]%
  {\begin{boitecours}[phantom={\signet{2}{#1}}]{coulProp}{#1}}{\end{boitecours}}
\newenvironment{methodebox}[1][Méthode]{\stepcounter{methode}%
  \begin{boitecours}[phantom={\signet[methode-\arabic{methode}]{2}{#1}}]{coulMeth}{#1}}%
  {\end{boitecours}}

\newcommand{\etape}[1]{\par\smallskip\textbf{\color{coulMeth}#1}\par\nobreak\smallskip}
\newcommand{\enonce}[1]{\emph{#1}\par\smallskip}
\newcommand{\reponse}[1]{\par\smallskip
  {\footnotesize\itshape\color{black!55}Réponses : #1}}
\newcommand{\demo}[1]{\textbf{\color{coulMeth}Démonstration.} #1}
\newcommand{\afaire}[1]{%
  \par\smallskip
  \noindent{\small\color{coulExo}$\blacktriangleright$\ \textbf{À faire :}
    fiche d'exercices du chapitre, #1.}\par\smallskip}

\newenvironment{calculs}{\setlength{\jot}{5pt}\@fleqntrue\@mathmargin\retraitcalculs
  \setlength{\abovedisplayskip}{4pt}\setlength{\belowdisplayskip}{4pt}%
  \setlength{\abovedisplayshortskip}{4pt}\setlength{\belowdisplayshortskip}{4pt}%
  \csname alignat*\endcsname{2}}%
  {\csname endalignat*\endcsname}
\newcommand{\rem}[1]{\text{\small\itshape\color{black!60}#1}}
\newcommand{\explique}[2]{\par\smallskip\noindent
  \begin{minipage}[t]{0.57\linewidth}\raggedright #1\end{minipage}\hfill
  \begin{minipage}[t]{0.40\linewidth}\raggedright\leavevmode
    {\small\itshape\color{black!60}#2}\end{minipage}\par}
\newcommand{\cas}[1]{\par\smallskip\noindent\textbf{\color{coulExo}#1}\nobreak}
\newcommand{\calc}[1]{\par\smallskip\textbf{\color{coulExo}#1}\par\nobreak}

\newtcolorbox{remarquebox}[1][Remarque]{%
  enhanced,breakable,colback=white,colframe=coulRem,
  boxrule=0pt,leftrule=2.5pt,arc=0pt,outer arc=0pt,
  left=8pt,right=6pt,top=5pt,bottom=5pt,
  before upper={\textbf{\color{coulRem}#1}\par\smallskip}}

\newtcolorbox{exemplebox}[1][Exemple]{%
  enhanced,breakable,colback=white,colframe=coulEx,
  boxrule=0pt,leftrule=2.5pt,arc=0pt,outer arc=0pt,
  left=8pt,right=6pt,top=5pt,bottom=5pt,
  before upper={\textbf{\color{coulEx}#1}\par\smallskip}}

\pgfplotsset{
  repere/.style={
    axis lines=middle,
    axis line style={-{Stealth[length=2.4mm]},thick},
    every axis x label/.style={at={(ticklabel* cs:1.0)},anchor=west},
    every axis y label/.style={at={(ticklabel* cs:1.0)},anchor=south},
    label style={font=\footnotesize},
    tick label style={font=\scriptsize},
    grid=both,
    major grid style={grille},
    minor tick num=0,
    tick align=inside,
    clip mode=individual,
  },
}
\tikzset{
  courbe/.style={coulCourbe,line width=1.1pt,smooth},
  courbeB/.style={coulCourbeB,line width=1.1pt,smooth},
  courbeC/.style={coulCourbeC,line width=1.1pt,smooth},
}

\newlist{etapes}{enumerate}{1}
\setlist[etapes]{label=\textbf{\arabic*.},leftmargin=*,itemsep=1pt,topsep=2pt}

\newlist{expressions}{enumerate}{1}
\setlist[expressions]{label={\Alph*\ $=$},leftmargin=*,itemsep=3pt,topsep=3pt}

\newcounter{exo}
\renewcommand{\theexo}{\ifnum\value{exo}<10 0\fi\arabic{exo}}
\newtcolorbox{exercice}[1][]{%
  enhanced,breakable,
  colback=white,colframe=coulExo,
  boxrule=0.9pt,arc=2pt,
  left=8pt,right=8pt,top=8pt,bottom=6pt,
  coltitle=white,fonttitle=\bfseries,
  attach boxed title to top left={xshift=8pt,yshift=-\tcboxedtitleheight/2},
  boxed title style={colback=coulExo,arc=2pt,boxrule=0pt},
  phantom={\signet{2}{Exercice \arabic{exo}}},
  title={Exercice \theexo\ifx\relax#1\relax\else{} --- #1\fi}}
\newenvironment{exo}[1][\relax]{\refstepcounter{exo}\begin{exercice}[#1]}{\end{exercice}}

  \newenvironment{correction}{\par}{\par}
  \newcommand{\autableau}{}
  \newcommand{\etapeexemples}{\etape{Exemples corrigés}}
  \newcommand{\etapetableau}[1]{\etape{#1}}
  \newenvironment{exempletableau}[1][Exemple]%
    {\begin{exemplebox}[#1]}{\end{exemplebox}}

\RequirePackage{comment}
\excludecomment{correctionavj}

\newcommand{\coursEntete}{}
\newcommand{\coursNiveau}{}
\newcommand{\coursTitre}{}
\newcommand{\coursSurTitre}{}
\newcommand{\coursSousTitre}{}

\newcommand{\enteteCours}[1]{\renewcommand{\coursEntete}{#1}}
\newcommand{\chapitrecours}[2]{\enteteCours{Chapitre #1 --- #2}}
\newcommand{\niveaucours}[1]{\renewcommand{\coursNiveau}{#1}}
\newcommand{\titrecours}[3][]{%
  \renewcommand{\coursTitre}{#2}%
  \renewcommand{\coursSurTitre}{#1}%
  \renewcommand{\coursSousTitre}{#3}}

\pagestyle{fancy}
\fancyhf{}
\fancyhead[L]{\footnotesize\color{coulPart}\textbf{\coursEntete}}
\fancyhead[R]{\footnotesize\color{coulPart}\textbf{\coursNiveau}}
\fancyfoot[C]{\footnotesize\thepage}
\renewcommand{\headrulewidth}{0.4pt}
\renewcommand{\headrule}{\hbox to\headwidth{\color{coulPart!40}\leaders\hrule height \headrulewidth\hfill}}

\setlength{\parskip}{2pt}

\AtBeginDocument{%
  \begin{center}
    \begin{tcolorbox}[enhanced,width=0.72\linewidth,colback=white,colframe=coulPart,
        boxrule=1.6pt,arc=3pt,halign=center,top=10pt,bottom=10pt]
      \ifx\coursSurTitre\empty
        {\Huge\bfseries\color{coulPart} \coursTitre}\\[4pt]
      \else
        {\Huge\bfseries\color{coulPart} \coursTitre}\\[3pt]
        {\Large\bfseries\color{coulPart} \coursSurTitre}\\[5pt]
      \fi
      {\normalsize \coursSousTitre}%
      
    \end{tcolorbox}
  \end{center}%
}
\makeatother

\chapitrecours{4}{Dérivation}
\niveaucours{1\textsuperscript{re} spécialité mathématiques}
\titrecours{DÉRIVATION}{Cours --- Première spécialité}

\begin{document}

\begin{remarquebox}[Comment utiliser ce cours]
Chaque partie commence par ce qu'il faut \textbf{savoir} (définitions, théorèmes), puis
vient ce qu'il faut \textbf{savoir faire} : une méthode, avec des
\textbf{exemples corrigés} faits en classe, puis des exercices
\og\textbf{À vous de jouer}\fg{} à chercher seul.

\end{remarquebox}

\partiecours{Limite en zéro d'une fonction}

\begin{exemplebox}[Exemple --- S'approcher d'une valeur interdite]
Certaines fonctions ne se calculent pas en $0$ : $0$ est une \textbf{valeur interdite}.
C'est le cas de
\[
  f(x)=\frac{x^{2}+2x}{x},
\]
car on ne peut pas diviser par $0$. On peut pourtant calculer $f(x)$ pour des valeurs de
$x$ de plus en plus proches de $0$ :
\begin{center}
\renewcommand{\arraystretch}{1.3}
\begin{tabular}{|c|c|c|c|c|c|c|c|}
\hline
$x$    & $-0{,}1$ & $-0{,}01$ & $-0{,}001$ & $0$ & $0{,}001$ & $0{,}01$ & $0{,}1$\\ \hline
$f(x)$ & $1{,}9$ & $1{,}99$ & $1{,}999$ & interdit & $2{,}001$ & $2{,}01$ & $2{,}1$\\ \hline
\end{tabular}
\end{center}
Quand $x$ s'approche de $0$, $f(x)$ s'approche de $2$.
\end{exemplebox}

\begin{definitionbox}[Définition 1 --- Limite en zéro]
Si $f(x)$ s'approche d'un nombre $L$ quand $x$ s'approche de $0$, on écrit
\[
  \lim_{x\to 0} f(x)=L
\]
et on lit « la limite de $f(x)$ quand $x$ tend vers $0$ est $L$ ». Dans l'exemple,
$\lim\limits_{x\to 0} f(x)=2$.
\end{definitionbox}

\begin{methodebox}[Méthode 1 --- Calculer une limite en zéro]
\etapeexemples
Calculer la limite en $0$ de chaque fonction.
\begin{questions}[label=\alph*)]
  \item $f(x)=\dfrac{x^{2}+2x}{x}$ (l'exemple ci-dessus)
  \item $g(x)=\dfrac{3x^{2}+5x}{x}$
  \item $h(x)=\dfrac{(2+x)^{2}-4}{x}$
\end{questions}
\begin{correction}
\cas{a)}
\begin{calculs}
f(x) &= \frac{x^{2}+2x}{x} &\qquad& \rem{on recopie l'expression}\\
f(x) &= \frac{\barre{x}(x+2)}{\barre{x}} &\qquad& \rem{on factorise le numérateur par $x$}\\
f(x) &= x+2 &\qquad& \rem{on simplifie : $x\ne 0$}\\
\lim_{x\to 0} f(x) &= 2 &\qquad& \rem{on remplace $x$ par $0$ dans $x+2$}
\end{calculs}
\cas{b)}
\begin{calculs}
g(x) &= \frac{3x^{2}+5x}{x} &&\\
g(x) &= \frac{\barre{x}(3x+5)}{\barre{x}} &&\\
g(x) &= 3x+5 &&\\
\lim_{x\to 0} g(x) &= 5 &\qquad& \rem{$3x$ tend vers $0$}
\end{calculs}
\cas{c)}
\begin{calculs}
h(x) &= \frac{(2+x)^{2}-4}{x} &&\\
h(x) &= \frac{4+4x+x^{2}-4}{x} &\qquad& \rem{identité remarquable}\\
h(x) &= \frac{\barre{x}(4+x)}{\barre{x}} &\qquad& \rem{$4$ et $-4$ s'annulent}\\
h(x) &= 4+x &&\\
\lim_{x\to 0} h(x) &= 4 &&
\end{calculs}
\end{correction}

\etape{À vous de jouer}
Calculer la limite en $0$ de chaque fonction :
\[
  k(x)=\frac{x^{2}+7x}{x} \qquad m(x)=\frac{5x^{2}-2x}{x} \qquad
  p(x)=\frac{(3+x)^{2}-9}{x}
\]
\reponse{$k(x)=x+7$, limite $7$ ;\; $m(x)=5x-2$, limite $-2$ ;\;
$p(x)=6+x$, limite $6$.}
\end{methodebox}

\afaire{exercices 1 à 4}

\partiecours{Nombre dérivé}

\begin{definitionbox}[Définition 2 --- Nombre dérivé]
\begin{minipage}[t]{0.56\linewidth}\raggedright
\vspace{0pt}
Soit $h\ne 0$. Le taux de variation de $f$ entre $a$ et $a+h$ est
\[
  T(h)=\frac{f(a+h)-f(a)}{h}.
\]
C'est le coefficient directeur de la sécante $(AM)$ (ch.~2).

\smallskip
Si $T(h)$ s'approche d'un nombre quand $h$ tend vers $0$ (partie I), on dit que $f$ est
\textbf{dérivable en $a$}. Cette limite est le \textbf{nombre dérivé} de $f$ en $a$,
noté $f'(a)$ :
\[
  f'(a)=\lim_{h\to 0}\frac{f(a+h)-f(a)}{h}.
\]
\end{minipage}\hfill
\begin{minipage}[t]{0.40\linewidth}\centering
\vspace{0pt}
\begin{tikzpicture}[x=0.95cm,y=0.95cm]
  \draw[-{Stealth[length=2mm]},thick] (-0.3,0) -- (4,0) node[below left,font=\scriptsize]{$x$};
  \draw[-{Stealth[length=2mm]},thick] (0,-0.3) -- (0,4) node[below left,font=\scriptsize]{$y$};
  \draw[coulCourbe,line width=1.1pt,samples=60,domain=0:3.3]
       plot (\x,{0.3*\x*\x+0.4});
  \draw[coulMeth,line width=0.9pt] (0.5,0.1) -- (3.55,3.76);
  \draw[coulExo,line width=0.9pt,dashed] (0,0.1) -- (3.6,2.26);
  \draw[dashed,thin] (1,0) -- (1,0.7);
  \draw[dashed,thin] (3,0) -- (3,3.1);
  \path (1,0.7) node[croix]{} node[above left,font=\scriptsize]{$A$};
  \path (3,3.1) node[croix]{} node[left,font=\scriptsize]{$M$};
  \node[below,font=\scriptsize] at (1,-0.03) {$a$};
  \node[below,font=\scriptsize] at (3,-0.03) {$a+h$};
  \node[coulMeth,font=\scriptsize,right] at (3.55,3.76) {sécante};
  \node[coulExo,font=\scriptsize,right] at (3.6,2.26) {tangente};
  \node[coulCourbe,font=\footnotesize] at (2.3,2.6) {$\Cf$};
\end{tikzpicture}

{\small Quand $h$ tend vers $0$, $M$ glisse vers $A$ : la sécante tourne et se
rapproche de la \textbf{tangente} en $A$.}
\end{minipage}
\end{definitionbox}

\begin{methodebox}[Méthode 2 --- Calculer un nombre dérivé par la définition]
\etapeexemples
\begin{questions}[label=\alph*)]
  \item $f(x)=x^{2}+3x$ : calculer $f'(1)$.
  \item $g(x)=\dfrac{1}{x}$ : calculer $g'(2)$.
\end{questions}
\begin{correction}
\cas{a)}
\begin{calculs}
f(1) &= 1+3 &\qquad& \rem{on calcule d'abord $f(a)$}\\
f(1) &= 4 &&\\
f(1+h) &= (1+h)^{2}+3(1+h) &\qquad& \rem{on remplace $x$ par $1+h$}\\
f(1+h) &= 1+2h+h^{2}+3+3h &&\\
f(1+h) &= h^{2}+5h+4 &&\\
T(h) &= \frac{h^{2}+5h+4-4}{h} &\qquad& \rem{$T(h)=\frac{f(1+h)-f(1)}{h}$}\\
T(h) &= \frac{\barre{h}(h+5)}{\barre{h}} &&\\
T(h) &= h+5 &&\\
f'(1) &= \lim_{h\to 0}\,(h+5) &\qquad& \rem{$f'(1)$ est la limite de $T(h)$ quand $h$ tend vers $0$}\\
f'(1) &= 5 &\qquad& \rem{on remplace $h$ par $0$ dans $h+5$}
\end{calculs}
\cas{b)}
\begin{calculs}
g(2+h)-g(2) &= \frac{1}{2+h}-\frac{1}{2} &\qquad& \rem{on calcule d'abord la différence}\\
g(2+h)-g(2) &= \frac{1\fois{2}}{(2+h)\fois{2}}-\frac{1\fois{(2+h)}}{2\fois{(2+h)}} &\qquad& \rem{même dénominateur}\\
g(2+h)-g(2) &= \frac{2-2-h}{2(2+h)} &&\\
g(2+h)-g(2) &= \frac{-h}{2(2+h)} &&\\
T(h) &= \frac{-\barre{h}}{2(2+h)}\times\frac{1}{\barre{h}} &\qquad& \rem{diviser par $h$, c'est multiplier par $\frac{1}{h}$}\\
T(h) &= \frac{-1}{2(2+h)} &&\\
g'(2) &= \lim_{h\to 0}\frac{-1}{2(2+h)} &&\\
g'(2) &= -\frac{1}{4} &\qquad& \rem{$2+h$ tend vers $2$}
\end{calculs}
\end{correction}

\etape{À vous de jouer}
\begin{questions}
  \item $h(x)=x^{2}-4x$ : calculer $h'(3)$.
  \item $k(x)=2x^{2}+1$ : calculer $k'(-1)$.
  \item $m(x)=\dfrac{1}{x}$ : calculer $m'(1)$.
\end{questions}
\reponse{1\textdegree) $T(h)=h+2$, $h'(3)=2$ ;\; 2\textdegree) $T(h)=2h-4$, $k'(-1)=-4$ ;\;
3\textdegree) $T(h)=\frac{-1}{1+h}$, $m'(1)=-1$.}
\end{methodebox}

\afaire{exercices 5 à 8}

\begin{remarquebox}[Une fonction qui n'est pas dérivable]
Pour la valeur absolue $f(x)=\lvert x\rvert$ en $0$ :
$T(h)=\dfrac{\lvert h\rvert}{h}$ vaut $1$ si $h>0$ et $-1$ si $h<0$. $T(h)$ ne se
rapproche pas d'un seul nombre : $f$ n'est \textbf{pas dérivable en $0$}. Sa courbe y
fait une pointe, sans tangente.
\end{remarquebox}

\partiecours{Tangente à une courbe}

\begin{definitionbox}[Définition 3 --- Tangente]
Soit $f$ dérivable en $a$ et $A\bigl(a\,;f(a)\bigr)$ le point de $\Cf$ d'abscisse $a$.
La \textbf{tangente} à $\Cf$ en $A$ est la droite qui passe par $A$ et dont le
coefficient directeur est $f'(a)$.
\end{definitionbox}

\begin{proprietebox}[Propriété 1 --- Équation de la tangente]
\begin{minipage}[c]{0.56\linewidth}\raggedright
La tangente à $\Cf$ au point d'abscisse $a$ a pour équation
\[
  y=f'(a)(x-a)+f(a).
\]
Elle passe par $A\bigl(a\,;f(a)\bigr)$ ; quand on avance de $1$, elle monte de $f'(a)$.
\end{minipage}\hfill
\begin{minipage}[c]{0.40\linewidth}\centering
\begin{tikzpicture}[x=0.85cm,y=0.85cm]
  \draw[-{Stealth[length=2mm]},thick] (-0.3,0) -- (4,0) node[below left,font=\scriptsize]{$x$};
  \draw[-{Stealth[length=2mm]},thick] (0,-0.3) -- (0,3.6) node[below left,font=\scriptsize]{$y$};
  \draw[coulCourbe,line width=1.1pt,samples=60,domain=0:3.5]
       plot (\x,{0.25*\x*\x+0.5});
  \draw[coulExo,line width=0.9pt] (0.5,0) -- (3.8,3.3);
  \draw[dashed,thin] (2,0) -- (2,1.5) -- (0,1.5);
  \draw[thin] (2,1.5) -- (3,1.5) -- (3,2.5);
  \node[below,font=\scriptsize] at (2.5,1.5) {$1$};
  \node[right,font=\scriptsize] at (3,2) {$f'(a)$};
  \path (2,1.5) node[croix]{} node[above left,font=\scriptsize]{$A$};
  \node[below,font=\scriptsize] at (2,-0.03) {$a$};
  \node[left,font=\scriptsize] at (-0.03,1.5) {$f(a)$};
  \node[coulCourbe,font=\footnotesize] at (1.1,1.4) {$\Cf$};
  \node[coulExo,font=\scriptsize,right] at (3.8,3.3) {tangente};
\end{tikzpicture}
\end{minipage}

\smallskip
\demo{La tangente a pour coefficient directeur $f'(a)$ et passe par $A$.}
\begin{calculs}
y &= f'(a)\,x+p &\qquad& \rem{équation d'une droite de coefficient directeur $f'(a)$}\\
f(a) &= f'(a)\times a+p &\qquad& \rem{$A\bigl(a\,;f(a)\bigr)$ est sur la droite}\\
p &= f(a)-f'(a)\times a &\qquad& \rem{on isole $p$}\\
y &= f'(a)\,x-f'(a)\times a+f(a) &\qquad& \rem{on remplace $p$ dans l'équation}\\
y &= f'(a)(x-a)+f(a) &\qquad& \rem{on factorise par $f'(a)$}
\end{calculs}
\end{proprietebox}

\begin{methodebox}[Méthode 3 --- Lire un nombre dérivé sur un graphique]
\begin{minipage}[c]{0.44\linewidth}\raggedright
On a tracé la courbe $\Cf$ d'une fonction $f$ et quatre de ses tangentes.

\etapeexemples
Lire $f'(-2)$ et $f'(0)$.

\etape{À vous de jouer}
Lire $f(-4)$, $f'(-4)$, puis $f(3)$ et $f'(3)$.
\reponse{$f(-4)=3$ et $f'(-4)=-3$ ;\; $f(3)=2$ et $f'(3)=-2$.}
\end{minipage}\hfill
\begin{minipage}[c]{0.54\linewidth}\centering
\begin{tikzpicture}[x=0.6cm,y=0.6cm]
  \repere{-5}{-2}{5}{5}
  \gradx{-4,-3,-2,1,2,3,4} \grady{-1,1,2,3,4}
  \draw[coulCourbe,line width=1.1pt]
    (-4,3) .. controls (-3.333,1) and (-2.667,-1) .. (-2,-1)
           .. controls (-1.333,-1) and (-0.667,-0.333) .. (0,1)
           .. controls (0.667,2.333) and (1.333,3) .. (2,3)
           .. controls (2.333,3) and (2.667,2.667) .. (3,2);
  \draw[coulExo,line width=0.9pt] (-4.333,4) -- (-2.667,-1);
  \draw[coulExo,line width=0.9pt] (-3.6,-1) -- (-0.4,-1);
  \draw[coulExo,line width=0.9pt] (-1.5,-2) -- (1.5,4);
  \draw[coulExo,line width=0.9pt] (2.25,3.5) -- (4.5,-1);
  \path[coulCourbe] (-4,3) node[croix]{} (-2,-1) node[croix]{} (0,1) node[croix]{}
    (3,2) node[croix]{};
  \node[coulCourbe,font=\footnotesize] at (2.4,3.6) {$\Cf$};
\end{tikzpicture}
\end{minipage}
\begin{correction}
\cas{$f'(-2)$}
\explique{La tangente en $-2$ est horizontale : $f'(-2)=0$.}{une droite horizontale a un
coefficient directeur nul}
\cas{$f'(0)$}
\explique{La tangente passe par $(0\,;1)$ et $(1\,;3)$ : quand on avance de $1$, elle monte
de $2$. Donc $f'(0)=2$.}{coefficient directeur $=\dfrac{\text{montée}}{\text{avancée}}$ ;
une descente compte en négatif}
\end{correction}
\end{methodebox}

\begin{methodebox}[Méthode 4 --- Déterminer une équation de la tangente]
\etapeexemples
\begin{questions}[label=\alph*)]
  \item $f(x)=x^{2}-2x+3$ : calculer $f'(2)$, puis une équation de la tangente à $\Cf$
        au point d'abscisse $2$.
  \item $g(x)=\dfrac{1}{x}$ : on a vu que $g'(2)=-\dfrac{1}{4}$ (méthode 2). Donner une
        équation de la tangente à $\Cg$ au point d'abscisse $2$.
\end{questions}
\begin{correction}
\cas{a)}
\begin{calculs}
f(2) &= 4-4+3 &&\\
f(2) &= 3 &&\\
f(2+h) &= (2+h)^{2}-2(2+h)+3 &&\\
f(2+h) &= h^{2}+2h+3 &\qquad& \rem{$4+4h+h^{2}-4-2h+3$}\\
T(h) &= \frac{h^{2}+2h}{h} &\qquad& \rem{$f(2+h)-f(2)$, divisé par $h$}\\
T(h) &= h+2 &&\\
f'(2) &= \lim_{h\to 0}\,(h+2) &&\\
f'(2) &= 2 &&\\
y &= f'(2)(x-2)+f(2) &\qquad& \rem{on recopie la formule}\\
y &= 2(x-2)+3 &&\\
y &= 2x-1 &&
\end{calculs}
\cas{b)}
\begin{calculs}
g(2) &= \frac{1}{2} &&\\
y &= g'(2)(x-2)+g(2) &&\\
y &= -\frac{1}{4}(x-2)+\frac{1}{2} &&\\
y &= -\frac{1}{4}x+1 &\qquad& \rem{$\frac{2}{4}+\frac{1}{2}=1$}
\end{calculs}
\end{correction}

\etape{À vous de jouer}
Avec les nombres dérivés déjà calculés à la méthode 2, donner une équation de la
tangente :
\begin{questions}
  \item à la courbe de $f(x)=x^{2}+3x$ au point d'abscisse $1$ ;
  \item à la courbe de $h(x)=x^{2}-4x$ au point d'abscisse $3$.
\end{questions}
\reponse{1\textdegree) $y=5(x-1)+4$, soit $y=5x-1$ ;\; 2\textdegree) $y=2(x-3)-3$, soit
$y=2x-9$.}
\end{methodebox}

\afaire{exercices 9 à 12}

\partiecours{Approximation linéaire}

\begin{proprietebox}[Propriété 2 --- Approximation linéaire]
\begin{minipage}[c]{0.52\linewidth}\raggedright
Si $f$ est dérivable en $a$, alors, pour $h$ proche de $0$ :
\[
  f(a+h)\approx f(a)+f'(a)\times h.
\]
Près de $A$, la courbe et sa tangente se confondent presque : on remplace l'ordonnée
$f(a+h)$ du point $M$ de la courbe par l'ordonnée $f(a)+f'(a)\times h$ du point $P$ de
la tangente. Plus $h$ est petit, plus l'écart $MP$ est petit.

C'est l'équation de la tangente (propriété 1), avec $x=a+h$.
\end{minipage}\hfill
\begin{minipage}[c]{0.44\linewidth}\centering
\begin{tikzpicture}[x=1.1cm,y=1.1cm]
  \draw[-{Stealth[length=2mm]},thick] (-0.3,0) -- (3.3,0) node[below left,font=\scriptsize]{$x$};
  \draw[-{Stealth[length=2mm]},thick] (0,-0.3) -- (0,3.2) node[below left,font=\scriptsize]{$y$};
  \draw[coulCourbe,line width=1.1pt,samples=60,domain=0:2.9]
       plot (\x,{0.3*\x*\x+0.4});
  \draw[coulExo,line width=0.9pt] (0,0.1) -- (3.1,1.96);
  \draw[dashed,thin] (1,0) -- (1,0.7);
  \draw[dashed,thin] (2.2,0) -- (2.2,1.852);
  \draw[dashed,thin] (2.2,1.42) -- (0,1.42);
  \draw[dashed,thin] (2.2,1.852) -- (0,1.852);
  \path (1,0.7) node[croix]{} node[above left,font=\scriptsize]{$A$};
  \path (2.2,1.852) node[croix]{} node[above left,font=\scriptsize]{$M$};
  \path (2.2,1.42) node[croix]{} node[below right,font=\scriptsize]{$P$};
  \node[below,font=\scriptsize] at (1,-0.03) {$a$};
  \node[below,font=\scriptsize] at (2.2,-0.03) {$a+h$};
  \node[left,font=\scriptsize] at (-0.03,1.92) {$f(a+h)$};
  \node[left,font=\scriptsize] at (-0.03,1.36) {$f(a)+f'(a)h$};
  \node[coulCourbe,font=\footnotesize] at (2.55,2.75) {$\Cf$};
  \node[coulExo,font=\scriptsize,right] at (3.1,1.96) {tangente};
\end{tikzpicture}
\end{minipage}
\end{proprietebox}

\begin{methodebox}[Méthode 5 --- Calculer une valeur approchée]
\etapeexemples
\begin{questions}[label=\alph*)]
  \item $f(x)=x^{2}$ : calculer $f'(5)$ par le taux d'accroissement, puis donner une
        valeur approchée de $5{,}03^{2}$ sans calculatrice.
  \item Un loyer augmente de $2$~\% par an pendant $3$ ans : il est multiplié par
        $1{,}02^{3}$. Avec $g(x)=x^{3}$, calculer $g'(1)$ par le taux d'accroissement,
        puis donner une valeur approchée de $1{,}02^{3}$ et le pourcentage
        d'augmentation sur les $3$ ans.
\end{questions}
\begin{correction}
\cas{a)}
\explique{Nombre dérivé en $5$ :}{méthode 2}
\begin{calculs}
f(5) &= 25 &&\\
f(5+h) &= 25+10h+h^{2} &\qquad& \rem{$(a+b)^{2}=a^{2}+2ab+b^{2}$}\\
T(h) &= \frac{10h+h^{2}}{h} &\qquad& \rem{$T(h)=\frac{f(5+h)-f(5)}{h}$}\\
T(h) &= 10+h &\qquad& \rem{on simplifie par $h$}\\
f'(5) &= 10 &\qquad& \rem{limite quand $h$ tend vers $0$}
\end{calculs}
\explique{Valeur approchée, avec $a=5$ et $h=0{,}03$ :}{$a$ proche, où $f(a)$ se calcule
de tête}
\begin{calculs}
f(5+0{,}03) &\approx f(5)+f'(5)\times 0{,}03 &\qquad& \rem{on recopie la formule}\\
5{,}03^{2} &\approx 25+10\times 0{,}03 &&\\
5{,}03^{2} &\approx 25{,}3 &\qquad& \rem{la calculatrice donne $25{,}3009$}
\end{calculs}
\cas{b)}
\explique{Nombre dérivé en $1$ :}{$(1+h)^{3}=(1+h)^{2}(1+h)$}
\begin{calculs}
g(1) &= 1 &&\\
g(1+h) &= (1+2h+h^{2})(1+h) &&\\
g(1+h) &= 1+3h+3h^{2}+h^{3} &\qquad& \rem{on développe}\\
T(h) &= \frac{3h+3h^{2}+h^{3}}{h} &&\\
T(h) &= 3+3h+h^{2} &\qquad& \rem{on simplifie par $h$}\\
g'(1) &= 3 &\qquad& \rem{limite quand $h$ tend vers $0$}
\end{calculs}
\explique{Valeur approchée, avec $a=1$ et $h=0{,}02$ :}{}
\begin{calculs}
g(1+0{,}02) &\approx g(1)+g'(1)\times 0{,}02 &\qquad& \rem{on recopie la formule}\\
1{,}02^{3} &\approx 1+3\times 0{,}02 &&\\
1{,}02^{3} &\approx 1{,}06 &\qquad& \rem{la calculatrice donne $1{,}061208$}
\end{calculs}
\explique{Le loyer augmente d'environ $6$~\% en $3$ ans.}{à peine plus que
$3\times 2$~\%}
\end{correction}

\etape{À vous de jouer}
\begin{questions}
  \item $f(x)=x^{2}$ : calculer $f'(4)$ par le taux d'accroissement, puis donner une
        valeur approchée de $4{,}01^{2}$ sans calculatrice.
  \item Un abonnement augmente de $4$~\% par an pendant $2$ ans : il est multiplié par
        $1{,}04^{2}$. Avec $g(x)=x^{2}$, calculer $g'(1)$ par le taux d'accroissement,
        puis donner une valeur approchée de $1{,}04^{2}$ et le pourcentage
        d'augmentation sur les $2$ ans.
\end{questions}
\reponse{1\textdegree) $f'(4)=8$ ; $4{,}01^{2}\approx 16+8\times 0{,}01=16{,}08$ ;\;
2\textdegree) $g'(1)=2$ ; $1{,}04^{2}\approx 1+2\times 0{,}04=1{,}08$ : environ
$+8$~\%.}
\end{methodebox}

\afaire{exercices 13 à 15}

\partiecours{Dérivées des fonctions polynômes}

\begin{definitionbox}[Définition 4 --- Fonction dérivée]
Si $f$ est dérivable en tout nombre $x$ d'un intervalle $I$, on dit que $f$ est
\textbf{dérivable sur $I$}. La fonction qui à $x$ associe $f'(x)$ est la
\textbf{fonction dérivée} de $f$, notée $f'$.

\smallskip
Elle \textbf{généralise le nombre dérivé} : au lieu de refaire le calcul de limite pour
chaque nombre $a$, on le fait une fois pour un $x$ quelconque. La fonction $f'$ donne
alors le nombre dérivé en \textbf{tout} $x$ de $I$ : il suffit de remplacer $x$ par $a$
pour obtenir $f'(a)$. Par exemple, si $f'(x)=2x$, alors $f'(3)=6$ et $f'(-1)=-2$.
\end{definitionbox}

\begin{proprietebox}[Propriété 3 --- Dérivée de la fonction carré]
Si $f(x)=x^{2}$, alors $f'(x)=2x$ pour tout réel $x$.

\smallskip
\demo{Pour un réel $a$ et $h\ne 0$ :}
\begin{calculs}
T(h) &= \frac{(a+h)^{2}-a^{2}}{h} &&\\
T(h) &= \frac{a^{2}+2ah+h^{2}-a^{2}}{h} &&\\
T(h) &= \frac{\barre{h}(2a+h)}{\barre{h}} &&\\
T(h) &= 2a+h &&\\
f'(a) &= \lim_{h\to 0}\,(2a+h) &&\\
f'(a) &= 2a &\qquad& \rem{$h$ tend vers $0$}
\end{calculs}
\end{proprietebox}

\begin{proprietebox}[Propriété 4 --- Dérivées des polynômes (admis)]
\begin{center}
\renewcommand{\arraystretch}{1.8}\setlength{\tabcolsep}{14pt}
\begin{tabular}{|c|c|}
\hline
$f(x)$ & $f'(x)$\\ \hline
$k$ (constante) & $0$\\ \hline
$mx+p$ & $m$\\ \hline
$x^{2}$ & $2x$\\ \hline
$x^{3}$ & $3x^{2}$\\ \hline
$x^{n}$ ($n\ge 1$ entier) & $nx^{n-1}$\\ \hline
\end{tabular}
\end{center}
Un polynôme se dérive \textbf{terme à terme}, et un coefficient se garde :
la dérivée de $kx^{n}$ est $k\times nx^{n-1}$.
\end{proprietebox}

\begin{methodebox}[Méthode 6 --- Dériver un polynôme]
\etapeexemples
Calculer la dérivée de chaque fonction.
\begin{questions}[label=\alph*)]
  \item $f(x)=4x^{3}-5x^{2}+2x-7$
  \item $g(x)=x^{5}-3x^{4}$
  \item $h(x)=\dfrac{1}{3}x^{3}+2x^{2}-x$
  \item $k(x)=(2x+1)^{2}$
\end{questions}
\begin{correction}
\cas{a)}
\begin{calculs}
f'(x) &= 4\times 3x^{2}-5\times 2x+2 &\qquad& \rem{$x^{3}$ donne $3x^{2}$ ; $2x$ donne $2$ ; $-7$ donne $0$}\\
f'(x) &= 12x^{2}-10x+2 &&
\end{calculs}
\cas{b)}
\begin{calculs}
g'(x) &= 5x^{4}-3\times 4x^{3} &\qquad& \rem{$x^{n}$ donne $nx^{n-1}$}\\
g'(x) &= 5x^{4}-12x^{3} &&
\end{calculs}
\cas{c)}
\begin{calculs}
h'(x) &= \frac{1}{3}\times 3x^{2}+2\times 2x-1 &&\\
h'(x) &= x^{2}+4x-1 &&
\end{calculs}
\cas{d)}
\begin{calculs}
k(x) &= 4x^{2}+4x+1 &\qquad& \rem{on développe d'abord}\\
k'(x) &= 8x+4 &&
\end{calculs}
\end{correction}

\etape{À vous de jouer}
Calculer la dérivée de chaque fonction :
\[
  f(x)=2x^{3}+x^{2}-6x+1 \qquad g(x)=-x^{4}+5x \qquad
  h(x)=\frac{1}{2}x^{4}-\frac{2}{3}x^{3} \qquad k(x)=(x-3)^{2}
\]
\reponse{$f'(x)=6x^{2}+2x-6$ ;\; $g'(x)=-4x^{3}+5$ ;\; $h'(x)=2x^{3}-2x^{2}$ ;\;
$k'(x)=2x-6$.}
\end{methodebox}

\afaire{exercices 16 à 20}

\partiecours{Dérivées des autres fonctions de référence}

\begin{proprietebox}[Propriété 5 --- Inverse, puissances inverses, racine carrée, cosinus, sinus]
\begin{center}
\renewcommand{\arraystretch}{2.4}\setlength{\tabcolsep}{14pt}
\begin{tabular}{|c|c|c|}
\hline
$f(x)$ & $f'(x)$ & $f$ est dérivable sur\\ \hline
$\dfrac{1}{x}$ & $-\dfrac{1}{x^{2}}$ & $\intoo{-\infty}{0}$ et $\intoo{0}{+\infty}$\\ \hline
$\dfrac{1}{x^{n}}$ ($n\ge 1$ entier) & $-\dfrac{n}{x^{n+1}}$ & $\intoo{-\infty}{0}$ et $\intoo{0}{+\infty}$\\ \hline
$\sqrt{x}$ & $\dfrac{1}{2\sqrt{x}}$ & $\intoo{0}{+\infty}$\\ \hline
$\cos x$ & $-\sin x$ & $\R$\\ \hline
$\sin x$ & $\cos x$ & $\R$\\ \hline
\end{tabular}
\end{center}
Les dérivées de $\dfrac{1}{x^{n}}$, de la racine carrée, du cosinus et du sinus sont
admises ; les fonctions cosinus et sinus seront étudiées au chapitre 9.
\end{proprietebox}

\begin{remarquebox}[Démonstration pour l'inverse]
On reprend la méthode 2 avec un nombre $a\ne 0$ quelconque. Pour $f(x)=\dfrac{1}{x}$ et
$h\ne 0$, on calcule d'abord la différence $f(a+h)-f(a)$, puis on divise par $h$.
\begin{calculs}
f(a+h)-f(a) &= \frac{1}{a+h}-\frac{1}{a} &\qquad& \rem{on recopie la différence}\\
f(a+h)-f(a) &= \frac{1\fois{a}}{(a+h)\fois{a}}-\frac{1\fois{(a+h)}}{a\fois{(a+h)}} &\qquad& \rem{même dénominateur}\\
f(a+h)-f(a) &= \frac{a-a-h}{a(a+h)} &&\\
f(a+h)-f(a) &= \frac{-h}{a(a+h)} &&\\
T(h) &= \frac{-\barre{h}}{a(a+h)}\times\frac{1}{\barre{h}} &\qquad& \rem{diviser par $h$, c'est multiplier par $\frac{1}{h}$}\\
T(h) &= \frac{-1}{a(a+h)} &&\\
f'(a) &= \lim_{h\to 0}\frac{-1}{a(a+h)} &&\\
f'(a) &= -\frac{1}{a^{2}} &\qquad& \rem{$a+h$ tend vers $a$}
\end{calculs}
\end{remarquebox}

\begin{methodebox}[Méthode 7 --- Utiliser les dérivées de référence]
\etapeexemples
\begin{questions}[label=\alph*)]
  \item $f(x)=\dfrac{1}{x}$ : calculer $f'(-2)$.
  \item $g(x)=\sqrt{x}$ : calculer $g'(9)$.
  \item Calculer la dérivée de $h(x)=\dfrac{5}{x}$, puis de $k(x)=4\sqrt{x}$.
  \item Calculer la dérivée de $m(x)=\dfrac{1}{x^{3}}$, puis de $p(x)=\dfrac{2}{x^{4}}$.
  \item Calculer la dérivée de $q(x)=3\sin x$, puis de $r(x)=-2\cos x$.
\end{questions}
\begin{correction}
\cas{a)}
\begin{calculs}
f'(x) &= -\frac{1}{x^{2}} &\qquad& \rem{on écrit d'abord la dérivée}\\
f'(-2) &= -\frac{1}{(-2)^{2}} &\qquad& \rem{puis on remplace $x$}\\
f'(-2) &= -\frac{1}{4} &&
\end{calculs}
\cas{b)}
\begin{calculs}
g'(x) &= \frac{1}{2\sqrt{x}} &&\\
g'(9) &= \frac{1}{2\times 3} &\qquad& \rem{$\sqrt{9}=3$}\\
g'(9) &= \frac{1}{6} &&
\end{calculs}
\cas{c)}
\begin{calculs}
h(x) &= 5\times\frac{1}{x} &\qquad& \rem{le coefficient $5$ se garde}\\
h'(x) &= 5\times\left(-\frac{1}{x^{2}}\right) &&\\
h'(x) &= -\frac{5}{x^{2}} &&\\
k'(x) &= 4\times\frac{1}{2\sqrt{x}} &&\\
k'(x) &= \frac{2}{\sqrt{x}} &\qquad& \rem{on simplifie par $2$}
\end{calculs}
\cas{d)}
\begin{calculs}
m'(x) &= -\frac{3}{x^{4}} &\qquad& \rem{$\frac{1}{x^{n}}$ donne $-\frac{n}{x^{n+1}}$, avec $n=3$}\\
p(x) &= 2\times\frac{1}{x^{4}} &\qquad& \rem{le coefficient $2$ se garde}\\
p'(x) &= 2\times\left(-\frac{4}{x^{5}}\right) &&\\
p'(x) &= -\frac{8}{x^{5}} &&
\end{calculs}
\cas{e)}
\begin{calculs}
q'(x) &= 3\cos x &\qquad& \rem{$\sin x$ donne $\cos x$}\\
r'(x) &= -2\times(-\sin x) &\qquad& \rem{$\cos x$ donne $-\sin x$}\\
r'(x) &= 2\sin x &\qquad& \rem{les deux signes « $-$ » donnent « $+$ »}
\end{calculs}
\end{correction}

\etape{À vous de jouer}
\begin{questions}
  \item $f(x)=\dfrac{1}{x}$ : calculer $f'(4)$. \quad $g(x)=\sqrt{x}$ : calculer $g'(25)$.
  \item Calculer la dérivée de chaque fonction :
        $h(x)=-\dfrac{2}{x}$ \quad $k(x)=6\sqrt{x}$ \quad $m(x)=\dfrac{1}{x^{2}}$ \quad
        $p(x)=\dfrac{3}{x^{5}}$.
  \item Calculer la dérivée de chaque fonction :
        $q(x)=4\cos x$ \quad $r(x)=-5\sin x$.
\end{questions}
\reponse{1\textdegree) $f'(4)=-\frac{1}{16}$ et $g'(25)=\frac{1}{10}$ ;\;
2\textdegree) $h'(x)=\frac{2}{x^{2}}$ ;\; $k'(x)=\frac{3}{\sqrt{x}}$ ;\;
$m'(x)=-\frac{2}{x^{3}}$ ;\; $p'(x)=-\frac{15}{x^{6}}$ ;\;
3\textdegree) $q'(x)=-4\sin x$ ;\; $r'(x)=-5\cos x$.}
\end{methodebox}

\afaire{exercices 21 à 24}

\partiecours{Dérivée d'une somme et d'un produit}

\begin{proprietebox}[Propriété 6 --- Somme, produit par un nombre, produit]
$u$ et $v$ sont deux fonctions dérivables sur un intervalle $I$, et $k$ un nombre réel.
\begin{center}
\renewcommand{\arraystretch}{1.8}\setlength{\tabcolsep}{14pt}
\begin{tabular}{|c|c|}
\hline
fonction & dérivée\\ \hline
$u+v$ & $u'+v'$\\ \hline
$ku$ & $ku'$\\ \hline
$uv$ & $u'v+uv'$\\ \hline
\end{tabular}
\end{center}
\textbf{Attention :} la dérivée d'un produit n'est \textbf{pas} le produit des dérivées.
\end{proprietebox}

\begin{remarquebox}[Démonstration pour le produit]
Pour $h\ne 0$, on enlève et on ajoute $u(a)\,v(a+h)$ au numérateur :
\begin{calculs}
T(h) &= \frac{u(a+h)\,v(a+h)-u(a)\,v(a)}{h} &&\\
T(h) &= \frac{u(a+h)\,v(a+h)-u(a)\,v(a+h)+u(a)\,v(a+h)-u(a)\,v(a)}{h} &&\\
T(h) &= \frac{u(a+h)-u(a)}{h}\times v(a+h)+u(a)\times\frac{v(a+h)-v(a)}{h} &&\\
(uv)'(a) &= \lim_{h\to 0} T(h) &&\\
(uv)'(a) &= u'(a)\,v(a)+u(a)\,v'(a) &&
\end{calculs}
En effet, quand $h$ tend vers $0$, les deux quotients tendent vers $u'(a)$ et $v'(a)$, et
$v(a+h)$ tend vers $v(a)$ (admis).
\end{remarquebox}

\begin{methodebox}[Méthode 8 --- Dériver une somme, un produit]
\etapeexemples
Calculer la dérivée de chaque fonction.
\begin{questions}[label=\alph*)]
  \item $f(x)=x^{2}+\dfrac{1}{x}$
  \item $g(x)=3x-4\sqrt{x}$
  \item $h(x)=(2x+1)(x^{2}-3)$
  \item $k(x)=x\sqrt{x}$ \;sur $\intoo{0}{+\infty}$
  \item $m(x)=x\cos x$
\end{questions}
\begin{correction}
\cas{a)}
\begin{calculs}
f'(x) &= 2x+\left(-\frac{1}{x^{2}}\right) &\qquad& \rem{somme : on dérive chaque terme}\\
f'(x) &= 2x-\frac{1}{x^{2}} &&
\end{calculs}
\cas{b)}
\begin{calculs}
g'(x) &= 3-4\times\frac{1}{2\sqrt{x}} &&\\
g'(x) &= 3-\frac{2}{\sqrt{x}} &&
\end{calculs}
\cas{c)}
\begin{calculs}
u(x) &= 2x+1 &\qquad& \rem{on nomme les deux facteurs}\\
u'(x) &= 2 &&\\
v(x) &= x^{2}-3 &&\\
v'(x) &= 2x &&\\
h'(x) &= 2(x^{2}-3)+(2x+1)\times 2x &\qquad& \rem{$u'v+uv'$}\\
h'(x) &= 2x^{2}-6+4x^{2}+2x &&\\
h'(x) &= 6x^{2}+2x-6 &&
\end{calculs}
\cas{d)}
\begin{calculs}
u(x) &= x &\qquad& \rem{$u'(x)=1$}\\
v(x) &= \sqrt{x} &\qquad& \rem{$v'(x)=\frac{1}{2\sqrt{x}}$}\\
k'(x) &= 1\times\sqrt{x}+x\times\frac{1}{2\sqrt{x}} &&\\
k'(x) &= \sqrt{x}+\frac{\sqrt{x}}{2} &\qquad& \rem{$\frac{x}{\sqrt{x}}=\sqrt{x}$, car $x=\sqrt{x}\times\sqrt{x}$}\\
k'(x) &= \frac{3}{2}\sqrt{x} &&
\end{calculs}
\cas{e)}
\begin{calculs}
u(x) &= x &\qquad& \rem{$u'(x)=1$}\\
v(x) &= \cos x &\qquad& \rem{$v'(x)=-\sin x$}\\
m'(x) &= 1\times\cos x+x\times(-\sin x) &\qquad& \rem{$u'v+uv'$}\\
m'(x) &= \cos x-x\sin x &&
\end{calculs}
\end{correction}

\etape{À vous de jouer}
Calculer la dérivée de chaque fonction :
\[
  f(x)=x^{3}-\frac{2}{x} \qquad g(x)=5x^{2}+2\sqrt{x} \qquad
  h(x)=(x-4)(3x+2)
\]
\[
  k(x)=(x^{2}+1)\sqrt{x} \qquad m(x)=x^{2}+\sin x \qquad p(x)=x\sin x
\]
\reponse{$f'(x)=3x^{2}+\frac{2}{x^{2}}$ ;\; $g'(x)=10x+\frac{1}{\sqrt{x}}$ ;\;
$h'(x)=(3x+2)+3(x-4)=6x-10$ ;\; $k'(x)=2x\sqrt{x}+\frac{x^{2}+1}{2\sqrt{x}}$ ;\;
$m'(x)=2x+\cos x$ ;\; $p'(x)=\sin x+x\cos x$.}
\end{methodebox}

\afaire{exercices 25 à 29}

\partiecours{Dérivée d'un inverse et d'un quotient}

\begin{proprietebox}[Propriété 7 --- Inverse et quotient]
$u$ et $v$ sont deux fonctions dérivables sur un intervalle $I$, et $v$ ne s'annule pas
sur $I$.
\begin{center}
\renewcommand{\arraystretch}{2.6}\setlength{\tabcolsep}{16pt}
\begin{tabular}{|c|c|}
\hline
fonction & dérivée\\ \hline
$\dfrac{1}{v}$ & $-\dfrac{v'}{v^{2}}$\\ \hline
$\dfrac{u}{v}$ & $\dfrac{u'v-uv'}{v^{2}}$\\ \hline
\end{tabular}
\end{center}
\end{proprietebox}

\begin{methodebox}[Méthode 9 --- Dériver un inverse, un quotient]
\etapeexemples
Calculer la dérivée de chaque fonction, là où elle est définie.
\begin{questions}[label=\alph*)]
  \item $f(x)=\dfrac{1}{4x-3}$
  \item $g(x)=\dfrac{1}{x^{2}+1}$
  \item $h(x)=\dfrac{2x+1}{x-3}$
  \item $k(x)=\dfrac{x^{2}}{x+1}$
\end{questions}
\begin{correction}
\cas{a)}
\begin{calculs}
v(x) &= 4x-3 &\qquad& \rem{$v'(x)=4$}\\
f'(x) &= -\frac{4}{(4x-3)^{2}} &\qquad& \rem{$-\frac{v'}{v^{2}}$}
\end{calculs}
\cas{b)}
\begin{calculs}
v(x) &= x^{2}+1 &\qquad& \rem{$v'(x)=2x$}\\
g'(x) &= -\frac{2x}{(x^{2}+1)^{2}} &&
\end{calculs}
\cas{c)}
\begin{calculs}
u(x) &= 2x+1 &\qquad& \rem{$u'(x)=2$}\\
v(x) &= x-3 &\qquad& \rem{$v'(x)=1$}\\
h'(x) &= \frac{2(x-3)-(2x+1)\times 1}{(x-3)^{2}} &\qquad& \rem{$\frac{u'v-uv'}{v^{2}}$ : le dénominateur reste au carré}\\
h'(x) &= \frac{2x-6-2x-1}{(x-3)^{2}} &\qquad& \rem{le « $-$ » change les signes de $uv'$}\\
h'(x) &= \frac{-7}{(x-3)^{2}} &&
\end{calculs}
\cas{d)}
\begin{calculs}
u(x) &= x^{2} &\qquad& \rem{$u'(x)=2x$}\\
v(x) &= x+1 &\qquad& \rem{$v'(x)=1$}\\
k'(x) &= \frac{2x(x+1)-x^{2}\times 1}{(x+1)^{2}} &&\\
k'(x) &= \frac{x^{2}+2x}{(x+1)^{2}} &\qquad& \rem{$2x^{2}+2x-x^{2}$}
\end{calculs}
\end{correction}

\etape{À vous de jouer}
Calculer la dérivée de chaque fonction :
\[
  f(x)=\frac{1}{5x+2} \qquad g(x)=\frac{1}{x^{2}-4x} \qquad
  h(x)=\frac{3x-1}{x+2} \qquad k(x)=\frac{x^{2}+3}{x}
\]
\reponse{$f'(x)=-\frac{5}{(5x+2)^{2}}$ ;\; $g'(x)=-\frac{2x-4}{(x^{2}-4x)^{2}}$ ;\;
$h'(x)=\frac{7}{(x+2)^{2}}$ ;\; $k'(x)=\frac{x^{2}-3}{x^{2}}$.}
\end{methodebox}

\afaire{exercices 30 à 35}

\besoinplace{12cm}
\partiecours{Dérivée des fonctions composées}

\begin{proprietebox}[Propriété 8 --- Dérivée d'une composée (admis)]
$u$ est une fonction dérivable sur un intervalle $I$, et $n\ge 2$ un entier.
\begin{center}
\renewcommand{\arraystretch}{2.4}\setlength{\tabcolsep}{16pt}
\begin{tabular}{|c|c|}
\hline
fonction & dérivée\\ \hline
$u^{2}$ & $2u'u$\\ \hline
$u^{n}$ & $nu'u^{n-1}$\\ \hline
$\dfrac{1}{u}$ & $-\dfrac{u'}{u^{2}}$\\ \hline
$\sqrt{u}$ & $\dfrac{u'}{2\sqrt{u}}$\\ \hline
$\cos u$ & $-u'\sin u$\\ \hline
$\sin u$ & $u'\cos u$\\ \hline
\end{tabular}
\end{center}
Pour $1/u$, $u$ ne s'annule pas sur $I$ (c'est la formule de la partie VIII) ; pour
$\sqrt{u}$, $u>0$ sur $I$.

Toutes ces lignes viennent d'une même formule : si $f(x)=g\big(u(x)\big)$, alors
\[
  f'(x)=u'(x)\times g'\big(u(x)\big).
\]
On dérive la fonction « extérieure » $g$ en gardant $u(x)$ à l'intérieur, puis on
multiplie par $u'(x)$.
\end{proprietebox}

\begin{remarquebox}[Démonstration pour $u^{2}$]
\begin{calculs}
(u^{2})' &= (u\times u)' &\qquad& \rem{$u^{2}$ est un produit}\\
(u^{2})' &= u'u+uu' &\qquad& \rem{$u'v+uv'$, avec $v=u$}\\
(u^{2})' &= 2u'u &&
\end{calculs}
\end{remarquebox}

\begin{methodebox}[Méthode 10 --- Dériver une fonction composée]
\etapeexemples
Calculer la dérivée de chaque fonction, là où elle est définie.
\begin{questions}[label=\alph*)]
  \item $f(x)=(2x+5)^{2}$
  \item $g(x)=(x^{2}+1)^{3}$
  \item $h(x)=\sqrt{3x+1}$
  \item $k(x)=\cos(2x)$
  \item $m(x)=\sin(x^{2}+1)$
\end{questions}
\begin{correction}
\cas{a)}
\begin{calculs}
u(x) &= 2x+5 &\qquad& \rem{$u'(x)=2$}\\
f'(x) &= 2\times 2\times(2x+5) &\qquad& \rem{$2u'u$}\\
f'(x) &= 8x+20 &&
\end{calculs}
\cas{b)}
\begin{calculs}
u(x) &= x^{2}+1 &\qquad& \rem{$u'(x)=2x$}\\
g'(x) &= 3\times 2x\times(x^{2}+1)^{2} &\qquad& \rem{$nu'u^{n-1}$, avec $n=3$}\\
g'(x) &= 6x(x^{2}+1)^{2} &&
\end{calculs}
\cas{c)}
\begin{calculs}
u(x) &= 3x+1 &\qquad& \rem{$u'(x)=3$ ; $u(x)>0$ pour $x>-\frac13$}\\
h'(x) &= \frac{3}{2\sqrt{3x+1}} &\qquad& \rem{$\frac{u'}{2\sqrt{u}}$}
\end{calculs}
\cas{d)}
\begin{calculs}
u(x) &= 2x &\qquad& \rem{$u'(x)=2$}\\
k'(x) &= -2\sin(2x) &\qquad& \rem{$-u'\sin u$}
\end{calculs}
\cas{e)}
\begin{calculs}
u(x) &= x^{2}+1 &\qquad& \rem{$u'(x)=2x$}\\
m'(x) &= 2x\cos(x^{2}+1) &\qquad& \rem{$u'\cos u$}
\end{calculs}
\end{correction}

\etape{À vous de jouer}
Calculer la dérivée de chaque fonction, là où elle est définie :
\[
  f(x)=(3x-2)^{2} \qquad g(x)=(x^{2}-2)^{4} \qquad h(x)=\sqrt{5x+4}
\]
\[
  k(x)=\sin(4x) \qquad m(x)=\cos(x^{2})
\]
\reponse{$f'(x)=18x-12$ ;\; $g'(x)=8x(x^{2}-2)^{3}$ ;\;
$h'(x)=\frac{5}{2\sqrt{5x+4}}$ ;\; $k'(x)=4\cos(4x)$ ;\; $m'(x)=-2x\sin(x^{2})$.}
\end{methodebox}

\afaire{exercices 36 à 38}

\partiecours{Étude des variations d'une fonction}

\begin{proprietebox}[Théorème 1 --- Signe de la dérivée et variations]
Soit $f$ dérivable sur un \textbf{intervalle} $I$.
\begin{itemize}
  \item Si $f'(x)\ge 0$ pour tout $x$ de $I$, alors $f$ est croissante sur $I$.
  \item Si $f'(x)\le 0$ pour tout $x$ de $I$, alors $f$ est décroissante sur $I$.
  \item Si $f'(x)=0$ pour tout $x$ de $I$, alors $f$ est constante sur $I$.
\end{itemize}
Si $f'(x)>0$ (sauf en quelques valeurs isolées où $f'(x)=0$), $f$ est
\textbf{strictement} croissante ; de même pour $f'(x)<0$.

\smallskip
Une tangente qui monte (coefficient directeur positif) suit une courbe qui monte.
\end{proprietebox}

\begin{methodebox}[Méthode 11 --- Étudier les variations d'une fonction]
\etapeexemples
Dresser le tableau de variations de chaque fonction.
\begin{questions}[label=\alph*)]
  \item $f(x)=x^{2}-6x+1$ sur $\R$
  \item $g(x)=x^{3}-3x^{2}-9x+2$ sur $\R$
  \item $h(x)=x+\dfrac{4}{x}$ sur $\intoo{0}{+\infty}$
\end{questions}
\begin{correction}
\cas{a)}
\begin{calculs}
f'(x) &= 2x-6 &\qquad& \rem{on dérive}\\
2x-6 &= 0 &\qquad& \rem{on cherche où $f'$ s'annule}\\
x &= 3 &&\\
f(3) &= 9-18+1 &&\\
f(3) &= -8 &&
\end{calculs}
\explique{$2x-6$ est négatif avant $3$, positif après.}{affine, coefficient $2>0$ :
$-$ puis $+$}
\begin{center}
\begin{tikzpicture}[scale=0.9]
  \tkzTabInit[lgt=2,espcl=2.4]{$x$/0.8,$f'(x)$/0.8,$f(x)$/1.6}{$-\infty$,$3$,$+\infty$}
  \tkzTabLine{,-,z,+,}
  \tkzTabVar{+/{},-/$-8$,+/{}}
\end{tikzpicture}
\end{center}
\cas{b)}
\begin{calculs}
g'(x) &= 3x^{2}-6x-9 &\qquad& \rem{on dérive}
\end{calculs}
\explique{Racines du trinôme $3x^{2}-6x-9$ :}{on cherche où $g'$ s'annule}
\begin{calculs}
\Delta &= (-6)^{2}-4\times 3\times(-9)=144 &\qquad& \rem{$\sqrt{144}=12$}\\
x_{1} &= \frac{6-12}{6}=-1 &&\\
x_{2} &= \frac{6+12}{6}=3 &&
\end{calculs}
\explique{Images des racines :}{les nombres de la ligne des $x$}
\begin{calculs}
g(-1) &= -1-3+9+2=7 &&\\
g(3) &= 27-27-27+2=-25 &&
\end{calculs}
\explique{$g'(x)$ est du signe de $a=3$, sauf entre les racines.}{signe du trinôme (ch.~3)}
\begin{center}
\begin{tikzpicture}[scale=0.9]
  \tkzTabInit[lgt=2,espcl=2.4]{$x$/0.8,$g'(x)$/0.8,$g(x)$/1.6}{$-\infty$,$-1$,$3$,$+\infty$}
  \tkzTabLine{,+,z,-,z,+,}
  \tkzTabVar{-/{},+/$7$,-/$-25$,+/{}}
\end{tikzpicture}
\end{center}
\cas{c)}
\begin{calculs}
h'(x) &= 1-\frac{4}{x^{2}} &\qquad& \rem{$4\times\left(-\frac{1}{x^{2}}\right)$}\\
h'(x) &= \frac{1\fois{x^{2}}}{1\fois{x^{2}}}-\frac{4}{x^{2}} &\qquad& \rem{même dénominateur}\\
h'(x) &= \frac{x^{2}-4}{x^{2}} &&\\
h'(x) &= \frac{(x-2)(x+2)}{x^{2}} &\qquad& \rem{$a^{2}-b^{2}=(a-b)(a+b)$}\\
h(2) &= 2+\frac{4}{2} &&\\
h(2) &= 4 &&
\end{calculs}
\explique{Sur $\intoo{0}{+\infty}$, $x^{2}>0$ et $x+2>0$ : $h'(x)$ a le signe de $x-2$.}{on
ne garde que les facteurs dont le signe change}
\begin{center}
\begin{tikzpicture}[scale=0.9]
  \tkzTabInit[lgt=2,espcl=2.4]{$x$/0.8,$h'(x)$/0.8,$h(x)$/1.6}{$0$,$2$,$+\infty$}
  \tkzTabLine{d,-,z,+,}
  \tkzTabVar{D+/{},-/$4$,+/{}}
\end{tikzpicture}
\end{center}
\end{correction}

\etape{À vous de jouer}
Dresser le tableau de variations de chaque fonction :
\[
  f(x)=-2x^{2}+8x-3 \qquad g(x)=x^{3}-12x+1 \qquad
  h(x)=\frac{2x-1}{x+1}\ \text{sur}\ \intoo{-1}{+\infty}
\]
\reponse{$f'(x)=-4x+8$ : $f$ croissante, puis décroissante, maximum $f(2)=5$ ;\;
$g'(x)=3x^{2}-12$ : $g$ croissante, décroissante, croissante, avec $g(-2)=17$ et $g(2)=-15$ ;\;
$h'(x)=\frac{3}{(x+1)^{2}}>0$ : $h$ croissante.}
\end{methodebox}

\afaire{exercices 39 à 42}

\partiecours{Extremums d'une fonction}

\begin{definitionbox}[Définition 5 --- Maximum, minimum]
Soit $f$ définie sur un intervalle $I$ et $c$ un nombre de $I$.
\begin{itemize}
  \item $f(c)$ est le \textbf{maximum} de $f$ sur $I$ si $f(x)\le f(c)$ pour tout $x$ de $I$.
  \item $f(c)$ est le \textbf{minimum} de $f$ sur $I$ si $f(x)\ge f(c)$ pour tout $x$ de $I$.
\end{itemize}
Un maximum ou un minimum est un \textbf{extremum}. Il est \textbf{local} s'il ne vaut que
sur un intervalle ouvert autour de $c$.
\end{definitionbox}

\begin{proprietebox}[Théorème 2 --- Extremum et dérivée]
Soit $f$ dérivable sur un intervalle ouvert $I$ et $c$ un nombre de $I$. Si $f'$
\textbf{s'annule en $c$ en changeant de signe}, alors $f$ a un \textbf{extremum local} en
$c$, et la tangente y est horizontale.

\smallskip
\begin{minipage}[t]{0.48\linewidth}\centering
{\small $f'$ passe de $+$ à $-$ : \textbf{maximum} $f(c)$}\\[4pt]
\begin{tikzpicture}[scale=0.85]
  \tkzTabInit[lgt=1.6,espcl=2]{$x$/0.8,$f'(x)$/0.8,$f(x)$/1.4}{,$c$,}
  \tkzTabLine{,+,z,-,}
  \tkzTabVar{-/{},+/$f(c)$,-/{}}
\end{tikzpicture}
\end{minipage}\hfill
\begin{minipage}[t]{0.48\linewidth}\centering
{\small $f'$ passe de $-$ à $+$ : \textbf{minimum} $f(c)$}\\[4pt]
\begin{tikzpicture}[scale=0.85]
  \tkzTabInit[lgt=1.6,espcl=2]{$x$/0.8,$f'(x)$/0.8,$f(x)$/1.4}{,$c$,}
  \tkzTabLine{,-,z,+,}
  \tkzTabVar{+/{},-/$f(c)$,+/{}}
\end{tikzpicture}
\end{minipage}
\end{proprietebox}

\begin{remarquebox}[Attention]
$f'(c)=0$ ne suffit pas : il faut que $f'$ \textbf{change de signe}. Pour $f(x)=x^{3}$,
$f'(x)=3x^{2}$ s'annule en $0$ mais reste positive : $f$ est croissante, sans extremum
en $0$.
\end{remarquebox}

\begin{methodebox}[Méthode 12 --- Déterminer les extremums d'une fonction]
\etapeexemples
$f(x)=x^{3}-3x+1$ sur $\intff{-2}{3}$.
\begin{questions}[label=\alph*)]
  \item Dresser le tableau de variations de $f$.
  \item Donner le maximum et le minimum de $f$ sur $\intff{-2}{3}$.
  \item Donner une équation de la tangente à $\Cf$ au point d'abscisse $-1$.
\end{questions}
\begin{correction}
\cas{a)}
\begin{calculs}
f'(x) &= 3x^{2}-3 &\qquad& \rem{on dérive}\\
f'(x) &= 3(x-1)(x+1) &\qquad& \rem{$3(x^{2}-1)$ : racines $-1$ et $1$}
\end{calculs}
\explique{Images des racines et des bornes :}{les nombres de la ligne des $x$}
\begin{calculs}
f(-2) &= -8+6+1=-1 &&\\
f(-1) &= -1+3+1=3 &&\\
f(1) &= 1-3+1=-1 &&\\
f(3) &= 27-9+1=19 &&
\end{calculs}
\begin{center}
\begin{tikzpicture}[scale=0.9]
  \tkzTabInit[lgt=2,espcl=2.2]{$x$/0.8,$f'(x)$/0.8,$f(x)$/1.6}{$-2$,$-1$,$1$,$3$}
  \tkzTabLine{,+,z,-,z,+,}
  \tkzTabVar{-/$-1$,+/$3$,-/$-1$,+/$19$}
\end{tikzpicture}
\end{center}
\cas{b)}
\explique{Le maximum de $f$ sur $\intff{-2}{3}$ est $19$, atteint en $3$.}{on compare les
« sommets » du tableau, bornes comprises}
\explique{Le minimum est $-1$, atteint en $-2$ et en $1$.}{$3$, atteint en $-1$, n'est qu'un
maximum local}
\cas{c)}
\explique{$f'(-1)=0$ : la tangente est horizontale, d'équation $y=3$.}{$y=f'(-1)(x+1)+f(-1)$
avec $f'(-1)=0$}
\end{correction}

\etape{À vous de jouer}
$g(x)=-x^{3}+3x^{2}$ sur $\intff{-1}{4}$.
\begin{questions}
  \item Dresser le tableau de variations de $g$.
  \item Donner le maximum et le minimum de $g$ sur $\intff{-1}{4}$.
  \item Donner une équation de la tangente à $\Cg$ au point d'abscisse $2$.
\end{questions}
\reponse{1\textdegree) $g'(x)=-3x(x-2)$ : « $-$ », « $+$ », « $-$ » ; $g(-1)=4$ et
$g(0)=0$ ; $g(2)=4$ et $g(4)=-16$ ;\; 2\textdegree) maximum $4$, atteint en $-1$ et en
$2$ ; minimum $-16$, atteint en $4$ ;\; 3\textdegree) $y=4$.}
\end{methodebox}

\afaire{exercices 43 à 45}

\partiecours{Étude complète d'une fonction}

\begin{remarquebox}[Plan d'étude d'une fonction]
\begin{etapes}
  \item \textbf{Ensemble de définition} : chercher les valeurs interdites (dénominateur
        nul, nombre négatif sous une racine).
  \item \textbf{Dérivée} : calculer $f'(x)$, puis l'écrire sous forme factorisée.
  \item \textbf{Signe de $f'(x)$} : étudier le signe de chaque facteur (tableau de signes).
  \item \textbf{Variations} : tableau de variations, avec les images des valeurs où $f'$
        s'annule et des bornes.
  \item \textbf{Extremums} : là où $f'(x)$ s'annule en changeant de signe.
  \item \textbf{Tangentes} : $y=f'(a)(x-a)+f(a)$ ; aux extremums, la tangente est
        horizontale.
  \item \textbf{Courbe} : tableau de valeurs, puis extremums, tangentes et courbe.
\end{etapes}
Dans un problème concret, la question posée dit quelles étapes faire ; on conclut par
une phrase, avec l'unité (partie suivante).
\end{remarquebox}

\begin{methodebox}[Méthode 13 --- Étudier une fonction de bout en bout]
\etapeexemples
On considère $f(x)=\dfrac{4x}{x^{2}+1}$.
\begin{questions}[label=\alph*)]
  \item Déterminer l'ensemble de définition de $f$.
  \item Calculer $f'(x)$ et l'écrire sous forme factorisée.
  \item Étudier le signe de $f'(x)$.
  \item Dresser le tableau de variations de $f$ et donner ses extremums.
  \item Donner une équation de la tangente $T$ à $\Cf$ au point d'abscisse $0$.
  \item Tracer $T$ et $\Cf$ sur $\intff{-4}{4}$.
\end{questions}
\begin{correction}
\cas{a)}
\explique{$x^{2}+1\ge 1$ : le dénominateur ne s'annule jamais.}{on cherche les valeurs
interdites}
\explique{$f$ est définie sur $\R$.}{aucune valeur interdite}
\cas{b)}
\begin{calculs}
u(x) &= 4x &\qquad& \rem{$u'(x)=4$}\\
v(x) &= x^{2}+1 &\qquad& \rem{$v'(x)=2x$}\\
f'(x) &= \frac{4(x^{2}+1)-4x\times 2x}{(x^{2}+1)^{2}} &\qquad& \rem{$\frac{u'v-uv'}{v^{2}}$}\\
f'(x) &= \frac{4-4x^{2}}{(x^{2}+1)^{2}} &\qquad& \rem{$4x^{2}+4-8x^{2}$}\\
f'(x) &= \frac{4(1-x)(1+x)}{(x^{2}+1)^{2}} &\qquad& \rem{$1-x^{2}=(1-x)(1+x)$}
\end{calculs}
\cas{c)}
\explique{$4>0$ et $(x^{2}+1)^{2}>0$ : $f'(x)$ a le signe de $1-x^{2}$.}{on ne garde que
les facteurs dont le signe change}
\explique{$1-x^{2}$ est négatif avant $-1$, positif entre $-1$ et $1$, négatif après.}{trinôme
de racines $\pm 1$, avec $a=-1<0$}
\cas{d)}
\begin{calculs}
f(-1) &= \frac{-4}{2}=-2 &\qquad& \rem{images des racines de $f'$}\\
f(1) &= \frac{4}{2}=2 &&
\end{calculs}
\begin{center}
\begin{tikzpicture}[scale=0.9]
  \tkzTabInit[lgt=2,espcl=2.2]{$x$/0.8,$f'(x)$/0.8,$f(x)$/1.6}{$-\infty$,$-1$,$1$,$+\infty$}
  \tkzTabLine{,-,z,+,z,-,}
  \tkzTabVar{+/{},-/$-2$,+/$2$,-/{}}
\end{tikzpicture}
\end{center}
\explique{$f$ a un minimum local $-2$ en $-1$ et un maximum local $2$ en $1$.}{$f'$ s'annule
en changeant de signe (théorème 2)}
\cas{e)}
\begin{calculs}
f(0) &= 0 &&\\
f'(0) &= \frac{4}{1}=4 &&\\
y &= 4(x-0)+0 &\qquad& \rem{$y=f'(0)(x-0)+f(0)$}\\
y &= 4x &&
\end{calculs}
\cas{f)}
\explique{Tableau de valeurs à la calculatrice (arrondis au dixième) :}{on ajoute les
extremums et le point de tangence}
\begin{center}
\renewcommand{\arraystretch}{1.3}
\begin{tabular}{|c|c|c|c|c|c|c|c|c|c|}
\hline
$x$ & $-4$ & $-3$ & $-2$ & $-1$ & $0$ & $1$ & $2$ & $3$ & $4$\\ \hline
$f(x)$ & $-0{,}9$ & $-1{,}2$ & $-1{,}6$ & $-2$ & $0$ & $2$ & $1{,}6$ & $1{,}2$ & $0{,}9$\\ \hline
\end{tabular}
\end{center}
\begin{center}
\begin{tikzpicture}[x=0.7cm,y=0.7cm]
  \repere{-5}{-3}{5}{3}
  \gradx{-4,-3,-2,-1,1,2,3,4} \grady{-2,-1,1,2}
  \draw[coulCourbe,line width=1.1pt,domain=-4:4,samples=120]
    plot (\x,{4*\x/(\x*\x+1)});
  \draw[coulExo,line width=0.9pt] (-0.65,-2.6) -- (0.65,2.6);
  \path[coulCourbe] (-1,-2) node[croix]{} (1,2) node[croix]{} (0,0) node[croix]{};
  \node[coulExo,font=\footnotesize] at (0.95,2.75) {$T$};
  \node[coulCourbe,font=\footnotesize] at (3.6,1.5) {$\Cf$};
\end{tikzpicture}
\end{center}
\explique{On place d'abord les extremums, où la tangente est horizontale, puis la tangente
$T$.}{la courbe « colle » à ses tangentes}
\end{correction}

\etape{À vous de jouer}
Mêmes questions a) à e) pour $g(x)=\dfrac{x^{2}+3}{x-1}$, avec la tangente au point
d'abscisse $0$ ; puis tracer $\Cg$ sur $\intff{-3}{0{,}5}$ et sur $\intff{1{,}5}{5}$.
\reponse{$g$ est définie sur $\R\setminus\{1\}$ ;
$g'(x)=\dfrac{(x+1)(x-3)}{(x-1)^{2}}$ ; maximum local $-2$ en $-1$, minimum local $6$ en
$3$ ; tangente en $0$ : $y=-3x-3$.}
\end{methodebox}

\afaire{exercices 46 à 48}

\partiecours{Problèmes concrets}

\begin{remarquebox}[Lire une fonction dans un contexte]
Quand $f(t)$ est une grandeur (altitude, concentration, nombre de personnes…) à
l'instant $t$ :
\begin{itemize}
  \item $f(0)$ est la valeur \textbf{au départ} ;
  \item $f'(t)$ est la \textbf{vitesse} à laquelle la grandeur varie à l'instant $t$
        (en km/h si $f(t)$ est en km et $t$ en heures) ;
  \item le maximum de $f$ donne la plus grande valeur, et l'instant où elle est atteinte.
\end{itemize}
Une durée en heures se convertit en minutes : $0{,}75$~h $=0{,}75\times 60$~min $=45$~min.
\end{remarquebox}

\begin{methodebox}[Méthode 14 --- Étudier une fonction dans un contexte]
\etapeexemples
On injecte un médicament à un patient à 8~h~45. La concentration du médicament dans le
sang, en mg/L, est donnée $t$ heures après l'injection par
\[
  C(t)=\frac{180t}{4t^{2}+9},\qquad t\in\intff{0}{8}.
\]
\begin{questions}[label=\alph*)]
  \item Calculer $C(0)$ et interpréter le résultat.
  \item Montrer que $C'(t)=\dfrac{180(9-4t^{2})}{(4t^{2}+9)^{2}}$.
  \item Étudier le signe de $C'(t)$ sur $\intff{0}{8}$, puis dresser le tableau de
        variations de $C$.
  \item À quelle heure la concentration est-elle maximale ? Que vaut-elle ?
  \item Le médicament agit quand la concentration est d'au moins $9$~mg/L. Pendant
        combien de temps agit-il ?
\end{questions}
\begin{correction}
\cas{a)}
\begin{calculs}
C(0) &= \frac{180\times 0}{4\times 0^{2}+9}=0 &&
\end{calculs}
\explique{Au moment de l'injection, il n'y a pas encore de médicament dans le sang.}{on
interprète dans le contexte}
\cas{b)}
\begin{calculs}
u(t) &= 180t &\qquad& \rem{$u'(t)=180$}\\
v(t) &= 4t^{2}+9 &\qquad& \rem{$v'(t)=8t$}\\
C'(t) &= \frac{180(4t^{2}+9)-180t\times 8t}{(4t^{2}+9)^{2}} &\qquad& \rem{$\frac{u'v-uv'}{v^{2}}$}\\
C'(t) &= \frac{720t^{2}+1620-1440t^{2}}{(4t^{2}+9)^{2}} &&\\
C'(t) &= \frac{1620-720t^{2}}{(4t^{2}+9)^{2}} &&\\
C'(t) &= \frac{180(9-4t^{2})}{(4t^{2}+9)^{2}} &\qquad& \rem{on factorise par $180$}
\end{calculs}
\cas{c)}
\explique{$180>0$ et $(4t^{2}+9)^{2}>0$ : $C'(t)$ a le signe de
$9-4t^{2}=(3-2t)(3+2t)$.}{on ne garde que les facteurs dont le signe change}
\explique{Sur $\intff{0}{8}$, $3+2t>0$ : $C'(t)$ a le signe de $3-2t$, qui s'annule en
$t=1{,}5$.}{$3-2t>0$ pour $t<1{,}5$}
\begin{calculs}
C(1{,}5) &= \frac{270}{9+9}=15 &\qquad& \rem{image de la racine de $C'$}\\
C(8) &= \frac{1440}{265}=\frac{288}{53} &\qquad& \rem{environ $5{,}4$}
\end{calculs}
\begin{center}
\begin{tikzpicture}[scale=0.9]
  \tkzTabInit[lgt=2,espcl=2.4]{$t$/0.8,$C'(t)$/0.8,$C(t)$/1.6}{$0$,$1{,}5$,$8$}
  \tkzTabLine{,+,z,-,}
  \tkzTabVar{-/$0$,+/$15$,-/$\frac{288}{53}$}
\end{tikzpicture}
\end{center}
\cas{d)}
\begin{calculs}
1{,}5~\text{h} &= 1~\text{h}+0{,}5\times 60~\text{min} &&\\
1{,}5~\text{h} &= 1~\text{h}~30~\text{min} &&
\end{calculs}
\explique{Le maximum est atteint $1$~h~$30$ après 8~h~45 : à \textbf{10~h~15}, la
concentration vaut \textbf{15~mg/L}.}{on répond par une phrase, avec l'unité}
\cas{e)}
\explique{On résout $C(t)\ge 9$ :}{on traduit la question}
\begin{calculs}
\frac{180t}{4t^{2}+9} &\ge 9 &&\\
180t &\ge 36t^{2}+81 &\qquad& \rem{on multiplie par $4t^{2}+9>0$}\\
36t^{2}-180t+81 &\le 0 &&\\
4t^{2}-20t+9 &\le 0 &\qquad& \rem{on divise par $9$}
\end{calculs}
\explique{Racines de $4t^{2}-20t+9$ :}{$\Delta=b^{2}-4ac$}
\begin{calculs}
\Delta &= 400-144=256 &&\\
t_{1} &= \frac{20-16}{8}=0{,}5 &&\\
t_{2} &= \frac{20+16}{8}=4{,}5 &&
\end{calculs}
\explique{Le trinôme est négatif entre ses racines : le médicament agit de $t=0{,}5$ à
$t=4{,}5$, soit pendant \textbf{4~heures} (de 9~h~15 à 13~h~15).}{signe de $a=4$ à
l'extérieur des racines}
\end{correction}

\etape{À vous de jouer}
Le nombre de spectateurs connectés à un match en ligne, en milliers, est donné $t$ heures
après 18~h par
\[
  N(t)=\frac{50t}{4t^{2}+25},\qquad t\in\intff{0}{4}.
\]
\begin{questions}
  \item Calculer $N(0)$ et interpréter le résultat.
  \item Montrer que $N'(t)=\dfrac{50(25-4t^{2})}{(4t^{2}+25)^{2}}$.
  \item Étudier le signe de $N'(t)$ sur $\intff{0}{4}$, puis dresser le tableau de
        variations de $N$.
  \item À quelle heure y a-t-il le plus de spectateurs ? Combien ?
  \item Le match est « en tendance » quand il y a au moins $2\,000$ spectateurs. Pendant
        combien de temps l'est-il ?
\end{questions}
\reponse{1\textdegree) $N(0)=0$ : personne au coup d'envoi ;\; 3\textdegree) $N'(t)$ a
le signe de $5-2t$ : $N$ croissante sur $\intff{0}{2{,}5}$, décroissante sur
$\intff{2{,}5}{4}$ ;\; 4\textdegree) $N(2{,}5)=2{,}5$ : $2\,500$ spectateurs à
20~h~30 ;\; 5\textdegree) $4t^{2}-25t+25\le 0$ pour $t\in\intff{1{,}25}{4}$ : pendant
2~h~45 (de 19~h~15 à 22~h).}
\end{methodebox}

\afaire{exercices 49 à 51}

\end{document}
